\chapter*{Introduction}
\addcontentsline{toc}{chapter}{Introduction}

\emph{This introduction is interim --- a brief orientation that stands in until the book's act-openers
are written. It is meant only to tell the reader, in advance, what kind of object the book builds and the
shape of the journey it makes.}

The construction this book describes begins from almost nothing: a single \emph{Fact} --- a proposition
together with a way to decide it, a yes-or-no with a procedure for telling which. From that one starting
point, and the bare act of telling two things apart, it assembles a finite apparatus: a measuring
instrument built by hand, one decidable step at a time, that never decrees a value but only reads what is
left over when a description is held against a baseline. Measurement, in these pages, is the reading of a
residue --- the part of the world a story fails to account for --- and the whole book is the slow
construction of an instrument honest enough to report that residue and nothing more. The instrument is
never asserted into being; it is shown, and it is shown through the historical experiments that recur in
every chapter --- Galileo's, Fourier's, Cantor's, and the rest --- each one the apparatus demonstrating a
power it has just earned rather than a decoration laid on after the fact.

The book is a single arc, a rise and a return, and it helps to carry the shape of it from the start. The
ascent begins at the bare difference --- the first act of telling two things apart --- and climbs,
faculty by faculty, until the apparatus can do the one thing it was not obviously able to do: turn and
read its own truth. At that turning point the distinction it began with collapses, true and false made to
coincide, and the book reaches its apex. What that collapse leaves behind --- the one quantity that
survives when every way of describing an action has been tried and divided away --- is the single
invariant the book is named for, and reaching it is the summit of the climb.

From the summit the reading reverses. The descent takes that one invariant and reads it back downward,
into matter: the first object it names is the electron, the bare difference of the opening pages returned
now with a charge, and what follows reads further physics --- a field, its radiation, the orientation of
spacetime --- off the same single number, each named with the same care. The return ends where the book
began, at the bare difference and the order it induces, now handed to the reader: the last apparatus in
the book, the one the construction cannot build and can only address. Every chapter closes by separating,
in plain words, what it has demonstrated from what it has assumed; the reader is asked to hold the
construction to exactly that standard, and no introduction, least of all an interim one, should ask for
more trust than the chapters themselves earn.
